\subsection{RQ4: How sequential is the decision problem?}
\label{sec:sequentiality}

Formally the problem is an MDP (\S\ref{sec:formulation}); the question is how
much a controller gains by looking beyond the current interval. We answer it
first in closed loop, then per decision.

\paragraph{Closed loop: clairvoyant lookahead as a controller.} Oracle-$H$
re-plans every interval with the true future: it rolls the cloned simulator
forward $H$ intervals from every legal action (the action, then no further
change) and takes the best. On the first five test seeds (35 episodes),
Oracle-1 reaches \LadderOraclehOne, against \LadderBandit{} for the bandit and
\LadderMilptrack{} for MILP-track. Extending the lookahead to three or six
intervals adds only \OracleThreeMinusOne{} [\OracleThreeMinusOneLo,
\OracleThreeMinusOneHi] and \OracleSixMinusOne{} [\OracleSixMinusOneLo,
\OracleSixMinusOneHi] (Oracle-6 is not better than Oracle-3). Even with
perfect foresight, lookahead beyond one interval is worth a few points to this
re-planning controller; knowing the next interval exactly is worth
\BanditBehindOracleOne{} over the bandit on the full test set. Oracle-1's lead
over the learners combines an exact model of the reward with knowledge of the
next interval's traffic, bursts and failures; we did not separate the two. For
this planner, what matters is the next interval, not the ones after it.

\paragraph{Per decision: open-loop rollouts.} On \SeqLiveStatesHOne{} states
visited by the greedy controller (diagnostic seeds 4001--4003, every fourth
decision; \SeqLiveStates{} of them have 24 remaining intervals), we compare the
action that is best for the first interval with the action that is best over
$H$ intervals under the same no-change continuation (Fig.~\ref{fig:seqdiag});
this is the effective-horizon idea of Laidlaw et
al.~\citep{laidlaw2023effective} with a no-change base policy. \emph{Under the
true future}, the first-interval choice is the 24-interval choice in only
\SeqLiveAgreeHTwentyFour\,\% of states and captures
\SeqLiveCapturedHTwentyFour\,\% of the 24-interval gain; in
\SeqLiveSacrificeHTwentyFour\,\% of states the 24-interval best move has a
negative immediate advantage. These open-loop values overstate what a
controller can realise: a move that pays off only after traffic changes can
also be made when the change arrives, which is what the re-planning oracle
does, and the closed-loop gains above are small. \emph{Holding traffic and
link state at their current values} during the rollouts, agreement is
\SeqFrozenAgreeHTwentyFour\,\%, the myopic choice captures
\SeqFrozenCapturedHTwentyFour\,\% of the gain and sacrifice moves occur in
\SeqFrozenSacrificeHTwentyFour\,\% of states, in every scenario
(Appendix~\ref{app:results}). With constant exogenous inputs and no further
moves, the $H$-interval value of a move is close to $H$ times its per-interval
gain minus its one-off cost, so the frozen agreement measures one thing: how
often amortizing a move's cost over a longer stay changes which move is best.
It rarely does. The open-loop gap between the two rollouts is anticipation of
exogenous change (traffic growth, bursts, failures and recoveries), which our
observation does not provide information for (no clock, no forecast; \S\ref{sec:rq3}
tests adding the clock). Neither rollout can detect value that arises from
\emph{later} moves---the opportunity cost of a dwell lock, or a move that
enables another under the protected-class rule---because the continuation
never moves again. The oracle ladder does not bound that value either: it
shows what a planner with the same no-change continuation gains, and a planner
with a better continuation could gain more. That Oracle-6 is not better than
Oracle-3 is consistent with this continuation bias.

\begin{figure}[t]
\centering
\includegraphics[width=\columnwidth]{fig_seqdiag}
\caption{Open-loop sequentiality diagnostic (greedy-visited states; from
\SeqLiveStatesHOne{} states at $H=1$ to \SeqLiveStates{} at $H=24$, since
states near an episode's end lack $H$ remaining intervals). Left: probability
that the action with the best first-interval reward is also the best over $H$
intervals (``take it, then change nothing''). Right: share of the best
$H$-interval gain that the myopic action captures. Solid: rollouts carry the
episode's true future traffic and failures (clairvoyant). Dashed: traffic and
link state held at their current values. Blue: base environment; orange:
delayed activation ($L=1$, \S\ref{sec:rq5}), where the myopic choice is
always no-op and captures none of the multi-interval gain.}
\label{fig:seqdiag}
\end{figure}
